\section{方法}
\label{sec:method}

\subsection{观察与问题形式化}
给定测试图像，冻结的 CLIP 视觉编码器输出 patch 特征 $\bF\in\R^{H\times W\times C}$ 及语义异常分数 $\Szero$（由 normal/abnormal 文本原型与 patch 特征的相似度得到）。理想情况下，异常区域应有较高的 $\Szero$；但由于 CLIP 低通，微小或纹理型缺陷的 $\Szero$ 往往并不显著。

另一方面，异常会扰动 $\bF$ 的高频分量。问题在于：高频响应同时包含异常与正常材质纹理，二者在幅度上不可区分。我们将这一现象总结为：

\begin{observation}\label{obs:ambiguity}
“这张图的正常高频水平”是与具体图像和材质相关（instance-specific）的量。零样本下没有训练数据可提前学得该量，因此任何图像无关的固定阈值都无法把异常从材质中分离出来。
\end{observation}

\subsection{读取高频}
我们在 CLIP patch 特征网格上做一层 Haar 离散小波变换（DWT）。低频 $\LL$ 对应物体的规则结构布局（正常应有的），高频子带给出局部突变的能量：
\begin{equation}
\HF = \tfrac{1}{C}\sum\nolimits_{c}\big(|\mathrm{LH}_c|+|\mathrm{HL}_c|+|\mathrm{HH}_c|\big).
\end{equation}
为把“物体自身的结构性边界”（正常的低频轮廓）与“无法被结构解释的突变”分开，我们定义边界感知可靠性
\begin{equation}
\W = \HF \cdot \big(1-\LFedge\big),\qquad \LFedge=\big\lVert\nabla \LL\big\rVert,
\end{equation}
其中 $\LFedge$ 为低频结构的梯度幅值。需强调：这一步只是\emph{读取} CLIP 已经编码的信息，不引入任何新信息，也不训练。

\subsection{证据选择}
我们用语义分数 $\Szero$ 与高频可靠性 $\W$（独立于 CLIP 语义）联合选择证据：
\begin{align}
\text{可信正常证据} &:\ \Szero\ \text{低} \ \wedge\ \W\ \text{低},\\
\text{异常证据} &:\ \Szero\ \text{高} \ \wedge\ \W\ \text{高}.
\end{align}
用 $\W$（而非仅用 $\Szero$）筛选可信正常区域是关键：若只用 CLIP 自身的语义置信，则筛出的“正常”里会混入 CLIP 看不见的高频异常，参照将被污染（自证闭环）。$\W$ 是 CLIP 语义之外的独立依据，才能把这类区域排除。

\subsection{逐图正常参照}
从选中的证据 patch 估计逐图视觉原型 $\vn,\va$，并对文本原型做保守更新：禁用异常原型更新（$\alpha_0=0$），仅对正常侧做小幅更新（小 $\beta_0$）。整个过程无训练、无反传、不更新 CLIP 参数；标签仅用于评测，不参与推理。最后以校准后的原型重算 patch 分数，异常图即高频信号相对该逐图参照的偏离。\figref{fig:arch} 给出整体框架。

\begin{figure*}[t]
\centering
\begin{tikzpicture}[
  font=\footnotesize,
  node distance=6mm,
  mod/.style={rounded corners=3pt,draw=black!35,thick,minimum height=1.05cm,align=center,inner sep=3pt},
  flow/.style={-{Latex[length=2mm]},thick,black!55},
]
% input
\node[mod,fill=cBase!30,minimum width=1.5cm] (img) {测试图\\\scriptsize 无标签};
% CLIP frozen
\node[mod,fill=cClip!30,right=of img,minimum width=1.7cm] (clip) {CLIP 编码器\\\scriptsize（冻结）};
% patch features
\node[mod,fill=cClip!15,right=of clip,minimum width=1.7cm] (feat) {Patch 特征\\\scriptsize $H\times W\times C$};
% semantic branch
\node[mod,fill=cClip!25,above right=2mm and 8mm of feat,minimum width=2.0cm] (sem) {语义分数\\\scriptsize normal/abnormal $\to S_0$};
% frequency branch
\node[mod,fill=cFreq!25,below right=2mm and 8mm of feat,minimum width=2.2cm] (freq) {读取高频\\\scriptsize Haar DWT: $\mathrm{HF}$\\\scriptsize $W=\mathrm{HF}\cdot(1-\mathrm{LF\text{-}edge})$};
% evidence
\node[mod,fill=cGlobal!25,right=14mm of feat,minimum width=2.0cm,yshift=0mm] (evi) {证据选择\\\scriptsize 正常: 低$S_0$ \& 低$W$\\\scriptsize 异常: 高$S_0$ \& 高$W$};
% per-image reference (highlight)
\node[mod,fill=cOurs!25,draw=cOurs!70,very thick,right=of evi,minimum width=2.2cm] (ref) {\textbf{逐图正常参照}\\\scriptsize 估计 $v_\mathrm{n},v_\mathrm{a}$\\\scriptsize 保守更新，无训练};
% output
\node[mod,fill=cBase!30,right=of ref,minimum width=1.6cm] (out) {异常图\\\scriptsize 相对参照的偏离};

\draw[flow] (img)--(clip);
\draw[flow] (clip)--(feat);
\draw[flow] (feat.east) to[out=30,in=180] (sem.west);
\draw[flow] (feat.east) to[out=-30,in=180] (freq.west);
\draw[flow] (sem.east) to[out=0,in=110] (evi.north);
\draw[flow] (freq.east) to[out=0,in=-110] (evi.south);
\draw[flow] (evi)--(ref);
\draw[flow] (ref)--(out);
\end{tikzpicture}
\caption{整体框架。测试图经冻结的 CLIP 得到 patch 特征，分为语义支（$S_0$）与频率支（Haar DWT $\to \mathrm{HF}/W$）；两者汇入证据选择，进而估计\emph{逐图正常参照}（保守更新，无训练），最后重算得到异常图。$W$ 提供独立于 CLIP 语义的高频依据。}
\label{fig:arch}
\end{figure*}


我们指出，multi-crop 聚合与 pixel-to-image 融合是标准的、与本文机制正交的增强手段，仅在\secref{sec:appendix}讨论，不计入核心方法。
